\documentclass[%
   % draft,     % Entwurfsstadium
   final,      % fertiges Dokument
%%%% --- Schriftgröße ---
   fontsize=12pt,
   %smallheadings,    % kleine Überschriften
   %normalheadings,   % normale Überschriften
   %bigheadings,       % große Überschriften
%%%% --- Sprache ---
%   ngerman,           % wird an andere Pakete weitergereicht
%%%% === Seitengröße ===
   % letterpaper,
   % legalpaper,
   % executivepaper,
   paper=a4,
   % a5paper,
   % landscap,
%%%% === Optionen für den Satzspiegel ===
   BCOR=5mm,          % Zusaetzlicher Rand auf der Innenseite
   %DIV11,            % Seitengroesse (siehe Koma Skript Dokumentation !)
   DIV=calc,         % automatische Berechnung einer guten Zeilenlaenge
   1.1headlines,     % Zeilenanzahl der Kopfzeilen
   %headinclude,     % Kopf einbeziehen
   %headexclude,      % Kopf nicht einbeziehen
   %footinclude,     % Fuss einbeziehen
   %footexclude,      % Fuss nicht einbeziehen
   %mpinclude,       % Margin einbeziehen
   %mpexclude,        % Margin nicht einbeziehen
   pagesize,         % Schreibt die Papiergroesse in die Datei.
                     % Wichtig für Konvertierungen
%%%% === Layout ===
   %oneside,         % einseitiges Layout
   twoside,          % Seitenraender für zweiseitiges Layout
   onecolumn,        % Einspaltig
   %twocolumn,       % Zweispaltig
   %openany,         % Kapitel beginnen auf jeder Seite
   open=right,        % Kapitel beginnen immer auf der rechten Seite
                     % (macht nur bei 'twoside' Sinn)
   %cleardoubleplain,    % leere, linke Seite mit Seitenstil 'plain'
   %cleardoubleempty,    % leere, linke Seite mit Seitenstil 'empty'
   titlepage,        % Titel als einzelne Seite ('titlepage' Umgebung)
   %notitlepage,     % Titel in Seite integriert
%%%% --- Absatzeinzug ---
   %                 % Absatzabstand: Einzeilig,
   %parskip,         % Freiraum in letzter Zeile: 1em
   %parskip*,        % Freiraum in letzter Zeile: Viertel einer Zeile
   %parskip+,        % Freiraum in letzter Zeile: Drittel einer Zeile
   %parskip-,        % Freiraum in letzter Zeile: keine Vorkehrungen
   %                 % Absatzabstand: Halbzeilig
   %halfparskip,     % Freiraum in letzter Zeile: 1em
   %halfparskip*,    % Freiraum in letzter Zeile: Viertel einer Zeile
   %halfparskip+,    % Freiraum in letzter Zeile: Drittel einer Zeile
   %halfparskip,     % Freiraum in letzter Zeile: keine Vorkehrungen
   %                 % Absatzabstand: keiner
   %parindent,        % Eingerückt (Standard)
%%%% --- Kolumnentitel ---
   headsepline,      % Linie unter Kolumnentitel
   %headnosepline,   % keine Linie unter Kolumnentitel
   %footsepline,     % Linie unter Fussnote
   %footnosepline,   % keine Linie unter Fussnote
%%%% --- Kapitel ---
   %chapterprefix,   % Ausgabe von 'Kapitel:'
   %nochapterprefix,  % keine Ausgabe von 'Kapitel:'
%%%% === Verzeichnisse (TOC, LOF, LOT, BIB) ===
   %liststotoc,      % Tabellen & Abbildungsverzeichnis ins TOC
   %idxtotoc,        % Index ins TOC
   %bibtotoc,         % Bibliographie ins TOC
   %bibtotocnumbered, % Bibliographie im TOC nummeriert
   %liststotocnumbered, % Alle Verzeichnisse im TOC nummeriert
   %tocindent,        % eingereuckte Gliederung
   %tocleft,         % Tabellenartige TOC
   %listsindent,      % eingereuckte LOT, LOF
   %listsleft,       % Tabellenartige LOT, LOF
   %pointednumbers,  % Überschriftnummerierung mit Punkt, siehe DUDEN !
   numbers=noenddot, % Überschriftnummerierung ohne Punkt, siehe DUDEN !
   %openbib,         % alternative Formatierung des Literaturverzeichnisses
%%%% === Matheformeln ===
   %leqno,           % Formelnummern links
   fleqn,            % Formeln werden linksbuendig angezeigt
]{scrbook}%     Klassen: scrartcl, scrreprt, scrbook

\makeatletter
\@ifpackagelater{scrbase}{2013/12/19}{%
  \KOMAoption{titlepage}{firstiscover}
}{
  \KOMAoption{titlepage}{yes}
}
\renewcommand*\toc@heading{%
 \chapter*{\contentsname}%
 \@mkboth{\MakeMarkcase{\contentsname}}{\MakeMarkcase{\contentsname}}%
}
\renewcommand*{\scr@float@listhead@warning}{%
}
\renewcommand*{\scr@float@addtolists@warning}{%
  \global\let\scr@float@addtolists@warning\relax
}
\makeatother
\listfiles
\usepackage{datenumber}
\usepackage[utf8x]{inputenc}
\usepackage{graphicx}


\def\svnId{$$Id:$$}
\makeindex

%% ------------------------------------------------------------------------
%%% Packages for LaTeX - programming
%
% Define commands that don't eat spaces.
\usepackage{xspace}
% IfThenElse
\usepackage{ifthen}
%%% Doc: ftp://tug.ctan.org/pub/tex-archive/macros/latex/contrib/oberdiek/ifpdf.sty
% command for testing for pdf-creation
\usepackage{ifpdf} %\ifpdf \else \fi

%%% Internal Commands: ----------------------------------------------
\makeatletter
%
\providecommand{\IfPackageLoaded}[2]{\@ifpackageloaded{#1}{#2}{}}
\providecommand{\IfPackageNotLoaded}[2]{\@ifpackageloaded{#1}{}{#2}}
\providecommand{\IfElsePackageLoaded}[3]{\@ifpackageloaded{#1}{#2}{#3}}
%
\newboolean{chapteravailable}%
\setboolean{chapteravailable}{false}%

\ifcsname chapter\endcsname
  \setboolean{chapteravailable}{true}%
\else
  \setboolean{chapteravailable}{false}%
\fi


\providecommand{\IfChapterDefined}[1]{\ifthenelse{\boolean{chapteravailable}}{#1}{}}%
\providecommand{\IfElseChapterDefined}[2]{\ifthenelse{\boolean{chapteravailable}}{#1}{#2}}%

\providecommand{\IfDefined}[2]{%
\ifcsname #1\endcsname
   #2 %
\else
     % do nothing
\fi
}

\providecommand{\IfElseDefined}[3]{%
\ifcsname #1\endcsname
   #2 %
\else
   #3 %
\fi
}

\providecommand{\IfElseUnDefined}[3]{%
\ifcsname #1\endcsname
   #3 %
\else
   #2 %
\fi
}


%
% Check for 'draft' mode - commands.
\newcommand{\IfNotDraft}[1]{\ifx\@draft\@undefined #1 \fi}
\newcommand{\IfNotDraftElse}[2]{\ifx\@draft\@undefined #1 \else #2 \fi}
\newcommand{\IfDraft}[1]{\ifx\@draft\@undefined \else #1 \fi}
%

% Definde frontmatter, mainmatter and backmatter if not defined
\@ifundefined{frontmatter}{%
   \newcommand{\frontmatter}{%
      %In Roemischen Buchstaben nummerieren (i, ii, iii)
      \pagenumbering{roman}
   }
}{}
\@ifundefined{mainmatter}{%
   % scrpage2 benoetigt den folgenden switch
   % wenn \mainmatter definiert ist.
   \newif\if@mainmatter\@mainmattertrue
   \newcommand{\mainmatter}{%
      % -- Seitennummerierung auf Arabische Zahlen zuruecksetzen (1,2,3)
      \pagenumbering{arabic}%
      \setcounter{page}{1}%
   }
}{}
\@ifundefined{backmatter}{%
   \newcommand{\backmatter}{
      %In Roemischen Buchstaben nummerieren (i, ii, iii)
      \pagenumbering{roman}
   }
}{}

% Pakete speichern die spaeter geladen werden sollen
\newcommand{\LoadPackagesNow}{}
\newcommand{\LoadPackageLater}[1]{%
   \g@addto@macro{\LoadPackagesNow}{%
      \usepackage{#1}%
   }%
}



\makeatother
%%% ----------------------------------------------------------------
%%
%%%% Doc: www.cs.brown.edu/system/software/latex/doc/calc.pdf
%% Calculation with LaTeX
%\usepackage{calc}
%
%%%% Doc: ftp://tug.ctan.org/pub/tex-archive/macros/latex/required/babel/babel.pdf
%% Languagesetting
%\usepackage[
%%	german,
%%	ngerman,
%	english,
%%	frensh,
%]{babel}

% \usepackage[%
%	final,
%	%draft % do not include images (faster)
% ]{graphicx}



\usepackage{marginnote}


%% Doc: (inside relsize.sty )
%% ftp://tug.ctan.org/pub/tex-archive/macros/latex/contrib/misc/relsize.sty
%  Set the font size relative to the current font size
%\usepackage{relsize}

%% Doc: ftp://tug.ctan.org/pub/tex-archive/macros/latex/contrib/ms/ragged2e.pdf
% Besserer Flatternsatz (Linksbuendig, statt Blocksatz)
%\usepackage{ragged2e}
%% Fonts
\usepackage{lmodern}


%% ~~~~~~~~~~~~~~~~~~~~~~~~~~~~~~~~~~~~~~~~~~~~~~~~~~~~~~~~~~~~~~~~~~~~~~~~
%% PDF related packages
%% ~~~~~~~~~~~~~~~~~~~~~~~~~~~~~~~~~~~~~~~~~~~~~~~~~~~~~~~~~~~~~~~~~~~~~~~~
%
%%% Doc: ftp://tug.ctan.org/pub/tex-archive/macros/latex/contrib/microtype/microtype.pdf
% Optischer Randausgleich mit pdfTeX
\ifpdf
\usepackage[%
	expansion=true, % better typography, but with much larger PDF file.
	protrusion=true
]{microtype}
\fi
%


%%% Doc: ftp://tug.ctan.org/pub/tex-archive/macros/latex/contrib/oberdiek/hypcap.pdf
% Links auf Gleitumgebungen springen nicht zur Beschriftung,
% sondern zum Anfang der Gleitumgebung
\IfPackageLoaded{hyperref}{%
	\usepackage[figure]{hypcap}
}



% Pakete Laden die nach Hyperref geladen werden sollen
\LoadPackagesNow % (ltxtable, tabularx)

% ~~~~~~~~~~~~~~~~~~~~~~~~~~~~~~~~~~~~~~~~~~~~~~~~~~~~~~~~~~~~~~~~~~~~~~~~
% figures and placement
% ~~~~~~~~~~~~~~~~~~~~~~~~~~~~~~~~~~~~~~~~~~~~~~~~~~~~~~~~~~~~~~~~~~~~~~~~

%% Bilder und Graphiken ==================================================

%%% Doc: only dtx Package
%\usepackage{float}             % Stellt die Option [H] f�r Floats zur Verf�gung

%%% Doc: No Documentation
%b\usepackage{flafter}          % Floats immer erst nach der Referenz setzen

% Defines a \FloatBarrier command, beyond which floats may not
% pass; useful, for example, to ensure all floats for a section
% appear before the next \section command.
%\usepackage[
%	section		% "\section" command will be redefined with "\FloatBarrier"
%]{placeins}
%
%%% Doc: ftp://tug.ctan.org/pub/tex-archive/macros/latex/contrib/subfig/subfig.pdf
% Incompatible: loads package capt-of. Loading of 'capt-of' afterwards will fail therefor
%\usepackage{subfig} % Layout wird weiter unten festgelegt !


% ~~~~~~~~~~~~~~~~~~~~~~~~~~~~~~~~~~~~~~~~~~~~~~~~~~~~~~~~~~~~~~~~~~~~~~~~
% verbatim packages
% ~~~~~~~~~~~~~~~~~~~~~~~~~~~~~~~~~~~~~~~~~~~~~~~~~~~~~~~~~~~~~~~~~~~~~~~~

%%% Doc: ftp://tug.ctan.org/pub/tex-archive/macros/latex/contrib/upquote/upquote.sty
\usepackage{upquote} % Setzt "richtige" Quotes in verbatim-Umgebung


%%% Doc: No documentation - documented in 'The LaTeX Companion'
% \usepackage{fancybox}   % for shadowbox, ovalbox

%%% Doc: ftp://tug.ctan.org/pub/tex-archive/macros/latex/contrib/misc/framed.sty
% \usepackage{framed}
% \renewcommand\FrameCommand{\fcolorbox{black}{shadecolor}}

\makeatletter
\IfPackageLoaded{framed}{%
   \IfPackageLoaded{marginnote}{%
������\begingroup
         \g@addto@macro\framed{%
���������   \let\marginnoteleftadjust\FrameSep
������������   \let\marginnoterightadjust\FrameSep
         }
      \makeatother
� }
}
\makeatother
%


\iffalse % CLASH mit dissertation.tex

% Layout mit 'geometry'
%%% Doc: ftp://tug.ctan.org/pub/tex-archive/macros/latex/contrib/geometry/manual.pdf
\usepackage{geometry}
\geometry{%
%%% Paper Groesse
   a4paper, % Andere a0paper, a1paper, a2paper, a3paper, , a5paper, a6paper,
            % b0paper, b1paper, b2paper, b3paper, b4paper, b5paper, b6paper
            % letterpaper, executivepaper, legalpaper
   %screen,  % a special paper size with (W,H) = (225mm,180mm)
   %paperwidth=,
   %paperheight=,
   %papersize=, %{ width , height }
   %landscape,  % Querformat
   portrait,    % Hochformat
%%% Koerper Groesse
   %hscale=,      % ratio of width of total body to \paperwidth
                  % hscale=0.8 is equivalent to width=0.8\paperwidth. (0.7 by default)
   %vscale=,      % ratio of height of total body to \paperheight
                  % vscale=0.9 is equivalent to height=0.9\paperheight.
   %scale=,       % ratio of total body to the paper. scale={ h-scale , v-scale }
   %totalwidth=,    % width of total body % (Generally, width >= textwidth)
   %totalheight=,   % height of total body, excluding header and footer by default
   %total=,        % total={ width , height }
   %textwidth=,    % modifies \textwidth, the width of body
   %textheight=,   % modifies \textheight, the height of body
   %body=,        % { width , height } sets both \textwidth and \textheight of the body of page.
   lines=44,       % enables users to specify \textheight by the number of lines.
   %includehead,  % includes the head of the page, \headheight and \headsep, into total body.
   %includefoot,  % includes the foot of the page, \footskip, into body.
   %includeheadfoot, % sets both includehead and includefoot to true
   %includemp,    % includes the margin notes, \marginparwidth and \marginparsep, into body
   %includeall,   % sets both includeheadfoot and includemp to true.
   %ignorehead,   % disregards the head of the page, headheight and headsep in determining vertical layout
   %ignorefoot,   % disregards the foot of page, footskip, in determining vertical layout
   %ignoreheadfoot, % sets both ignorehead and ignorefoot to true.
   %ignoremp,     % disregards the marginal notes in determining the horizontal margins
   ignoreall,     % sets both ignoreheadfoot and ignoremp to true
   heightrounded, % This option rounds \textheight to n-times (n: an integer) of \baselineskip
   %hdivide=,     % { left margin , width , right margin }
                  % Note that you should not specify all of the three parameters
   %vdivide=,     % { top margin , height , bottom margin }
   %divide=,      % ={A,B,C} %  is interpreted as hdivide={A,B,C} and vdivide={A,B,C}.
%%% Margin
   %left=,        % left margin (for oneside) or inner margin (for twoside) of total body
                  % alias: lmargin, inner
   %right=,       % right or outer margin of total body
                  % alias: rmargin outer
   %top=,         % top margin of the page.
                  % Alias : tmargin
   %bottom=,      % bottom margin of the page
                  % Alias : bmargin
   %hmargin=,     % left and right margin. hmargin={ left margin , right margin }
   %vmargin=,     % top and bottom margin. vmargin={ top margin , bottom margin }
   %margin=,      % margin={A,B} is equivalent to hmargin={A,B} and vmargin={A,B}
   %hmarginratio, % horizontal margin ratio of left (inner) to right (outer).
   %vmarginratio, % vertical margin ratio of top to bottom.
   %marginratio,  % marginratio={ horizontal ratio , vertical ratio }
   %hcentering,   % sets auto-centering horizontally and is equivalent to hmarginratio=1:1
   %vcentering,   % sets auto-centering vertically and is equivalent to vmarginratio=1:1
   centering,    % sets auto-centering and is equivalent to marginratio=1:1
   twoside,       % switches on twoside mode with left and right margins swapped on verso pages.
   %asymmetric,   % implements a twosided layout in which margins are not swapped on alternate pages
                  % and in which the marginal notes stay always on the same side.
   bindingoffset=5mm,  % removes a specified space for binding
%%% Dimensionen
   %headheight=,  % Alias:  head
   %headsep=,     % separation between header and text
   %footskip=,    % distance separation between baseline of last line of text and baseline of footer
                  % Alias: foot
   %nohead,       % eliminates spaces for the head of the page
                  % equivalent to both \headheight=0pt and \headsep=0pt.
   %nofoot,       % eliminates spaces for the foot of the page
                  % equivalent to \footskip=0pt.
   %noheadfoot,   % equivalent to nohead and nofoot.
   %footnotesep=, % changes the dimension \skip\footins,.
                  % separation between the bottom of text body and the top of footnote text
   marginparwidth=0pt, % width of the marginal notes
                  % Alias: marginpar
   %marginparsep=,% separation between body and marginal notes.
   %nomarginpar,  % shrinks spaces for marginal notes to 0pt
   %columnsep=,   % the separation between two columns in twocolumn mode.
   %hoffset=,
   %voffset=,
   %offset=,      % horizontal and vertical offset.
                  % offset={ hoffset , voffset }
   %twocolumn,    % twocolumn=false denotes onecolumn
   %twoside,
   textwidth=400pt,   % sets \textwidth directly
   %textheight=600pt,  % sets \textheight directly
   %reversemp,    % makes the marginal notes appear in the left (inner) margin
                  % Alias: reversemarginpar
} % Endif

\fi


\raggedbottom     % Variable Seitenhoehen zulassen



%% Schriften (Sections )==================================================

\IfElsePackageLoaded{fourier}{
   \newcommand\SectionFontStyle{\rmfamily}
}{
   \newcommand\SectionFontStyle{\sffamily}
}

% -- Koma Schriften --
\IfChapterDefined{%
   \setkomafont{chapter}{\huge\SectionFontStyle}    % Chapter
}
\setkomafont{descriptionlabel}{\SectionFontStyle}
\setkomafont{sectioning}{\SectionFontStyle} %  % Titelzeilen % \bfseries
\setkomafont{pagenumber}{\bfseries\SectionFontStyle}             % Seitenzahl
\setkomafont{pageheadfoot}{\small\sffamily}        % Kopfzeile \setkomafont{pagehead}{\small\sffamily}
%\setkomafont{pagefoot}{\small\sffamily}        % Kopfzeile
%\setkomafont{descriptionlabel}{\itshape}        % Kopfzeile
%

\renewcommand*{\raggedsection}{\raggedright} % Titelzeile linksbuendig, haengend
%
%% UeberSchriften (Chapter und Sections) =================================
% -- Ueberschriften komlett Umdefinieren --
%%% Doc: ftp://tug.ctan.org/pub/tex-archive/macros/latex/contrib/titlesec/titlesec.pdf

%\usepackage{titlesec} 

%\titleformat{\chapter}[block]
%{\vspace{-5pc}\usekomafont{chapter}\Large \color{black}}
%{\Huge \textsubscript{\thechapter} \filright }
%{1px}
%{\titlerule\vspace{.5pc}\\ \LARGE  }   % {before}[after] (before chaptertitle and after)
%  [\color{black} \vspace{0.6pc} \filright {\titlerule}]


%% Captions (Schrift, Aussehen) ==========================================

% % Folgende Befehle werden durch das Paket caption und subfig ersetzt !
% \setcapindent{1em} % Einrueckung der Beschriftung
% \setkomafont{caption}{\color{black}\small\sffamily\RaggedRight}  % Schrift f�r Caption
% \setkomafont{captionlabel}{\color{black}\small}   % Schrift f�r 'Abbildung' usw.

%%% Doc: ftp://tug.ctan.org/pub/tex-archive/macros/latex/contrib/caption/caption.pdf
\usepackage{caption}
% Aussehen der Captions
\captionsetup{
   margin = 10pt,
   font = {small,rm},
   labelfont = {small,bf},
   format = plain, % oder 'hang'
   indention = 0em,  % Einruecken der Beschriftung
   labelsep = colon, %period, space, quad, newline
   justification = RaggedRight, % justified, centering
   singlelinecheck = true, % false (true=bei einer Zeile immer zentrieren)
   position = bottom %top
}
%%% Bugfix Workaround
\DeclareCaptionOption{parskip}[]{}
\DeclareCaptionOption{parindent}[]{}

% Aussehen der Captions f�r subfigures (subfig-Paket)
\IfPackageLoaded{subfig}{
 \captionsetup[subfloat]{%
   margin = 10pt,
   font = {small,rm},
   labelfont = {small,bf},
   format = plain, % oder 'hang'
   indention = 0em,  % Einruecken der Beschriftung
   labelsep = space, %period, space, quad, newline
   justification = RaggedRight, % justified, centering
   singlelinecheck = true, % false (true=bei einer Zeile immer zentrieren)
   position = bottom, %top
   labelformat = parens % simple, empty % Wie die Bezeichnung gesetzt wird
 }
}

\usepackage{rotating}


% Malte
\newcommand{\bv}{behavioral variations\xspace}
\newcommand{\Bv}{Behavioral variations\xspace}
\newcommand{\BV}{Behavioral Variations\xspace}

\newcommand{\forsubject}{\code{subject}\xspace}
\newcommand{\context}{\code{context}\xspace}
\newcommand{\LPAREN}{\code{\textbf{(}}}
\newcommand{\RPAREN}{\code{\textbf{)}}}
\newcommand{\LBRACE}{\code{\textbf{\{}}}
\newcommand{\RBRACE}{\code{\textbf{\}}}}
\newcommand{\LBRACKET}{\code{\textbf{[}}}
\newcommand{\RBRACKET}{\code{\textbf{]}}}
%\newcommand{\ind}{\hspace*{1.5em}}
\newcommand{\nl}{\newline}
\newcommand{\OR}{~\textbar}
\newcommand{\LBr}{ \texttt{(} }
\newcommand{\RBr}{\texttt{)}}
\newcommand{\LB}{\texttt{\{~}}
\newcommand{\RB}{\texttt{\}}}
%\newcommand{\OR}{\textbar~}
\newcommand{\SEMICOLON}{\texttt{;}}

\newcommand{\supervisor}[2] { %
 %\begin{tabular}{l l}
 \par{%
   \begin{minipage}[t]{3cm}
     \large #1 :
   \end{minipage}\begin{minipage}[t]{10cm}
     \large #2 
   \end{minipage}
 }
 }




 %TEX root = main.tex
%%%%%%%%%%%%%%%%%%%%%%%%%%%%%%%%%%%%%%%%%%%%%%%%%%%%%%%%%%%%%%%%%%%%%%%%%%%%%%%%%%%%%%%%%%%%%%%%%%%
%
% This file contains:
%
% Comment Mode Support 
% ==================
% The following features are only active while in comment mode. To enter this mode, simply 
% set \commentmode in the preamble. 
%
% \documentclass{...}
% \input{auxiliary}
% \commentmode
%
% - Draft Notice and Svn Info in a Framed Header on Every Page. This feature is automatically 
%   activated while in comment mode, but needs an additional declaration in the preamble to set 
%   the SVN Id: \def\svnId{$$Id:$$}
% 
%   \documentclass{...}
%   \def\svnId{$$Id:$$}
%   \input{auxiliary}
%   \commentmode
%
% - Comments Labeled with Author Id. Define your own author id before \begin{document}:
%   \newcommand\myAuthorID[1]{\nbnote{MAID}{#1}}
%   Use it anywhere in the text: 
%   ... \myAuthorID{this is a comment} ... 
%
% Listings Package Configuration for Several Languages
% ====================================================
% - Syntax highlighting for JCop, ContextJ, LogicAJ, LogicAJ2, GenTL, Prolog
% - Code makro to define code snippets independently of their appeariance. Use:
%   \code{public void foo(){}}
%
% Published Notice on Title Page
% ============================= 
% Paper reference feature to place a publication note on top of the title page of local downloads. 
% Use the following command after \begin{document}: 
% \paperref{<framed box width>}{<height>}{This paper appeared at the n-th XY Conference on Z}
%
% last modified: MA, 2009-08-17_14-47
%%%%%%%%%%%%%%%%%%%%%%%%%%%%%%%%%%%%%%%%%%%%%%%%%%%%%%%%%%%%%%%%%%%%%%%%%%%%%%%%%%%%%%%%%%%%%%%%%%%
\usepackage{amsmath}
\usepackage[usenames,dvipsnames]{color}

\usepackage{framed}
\usepackage{xspace}
\usepackage[hyphens]{url}
%\urlstyle{same}
\urlstyle{tt}
%\usepackage{caption}
\DeclareMathAlphabet{\mathpzc}{OT1}{pzc}{m}{it}
% General %%%%%%%%%%%%%%%%%%%%%%%%%%%%%%%%%%%%%%%%%%%%%%%%%%%%%%%%%%%%%%%%%%%%%%%%%%%%%%%%%%%%%%%%%

% API Table %%%%%%%%%%%%%%%%%%%%%%%%%%%%%%%%%%%%%%%%%%%%%%%%%%%%%%%%%%%%%%%%%%%%%%%%%%%%%%%
\usepackage{longtable}
\usepackage{multirow}

\newenvironment{api2}[1]{%
\begin{tabular}{p{.5cm} p{13.5cm}}%
\hline%
%\multicolumn{2}{|c|}{\code{#1}}\\ \hline %
}%
{\hline%
\end{tabular}%
}
\newcommand{\apidesctxt}[2]{\multicolumn{2}{l}{\lstinline[language=JCop, basicstyle=\listingsfont]{public #1}} \\ & #2 \\}

\newenvironment{api}[1]{%
\begin{longtable}{p{.5cm} p{12.5cm}}%
%\hline%
%\multicolumn{2}{|c|}{\code{#1}}\\ \hline %
}%
{%\hline%
\end{longtable}%
}

\newcommand{\apidesc}[2]{\multicolumn{2}{l}{\lstinline[language=JCop, basicstyle=\listingsfont]{public #1}} \\ & \footnotesize{#2} \\}


% COMMENTS & DRAFT  %%%%%%%%%%%%%%%%%%%%%%%%%%%%%%%%%%%%%%%%%%%%%%%%%%%%%%%%%%%%%%%%%%%%%%%%%%%%%%%

\usepackage{ednotes}
\usepackage{amssymb}
\usepackage{datetime}
\usepackage{everypage}

% date and timeformat
\newdateformat{simpledate}{\THEYEAR-\twodigit{\THEMONTH}-\twodigit{\THEDAY}}
\newtimeformat{simpletime}{\twodigit{\THEHOUR}:\twodigit{\THEMINUTE}}
\newcommand{\timestamp}{\simpledate\today~\simpletime}

% header for draft notice
\newcommand{\draftheader}{%
  %\framebox[\textwidth][t]{%
	  \centerline{%
		  \textbf{DRAFT}
		  \scriptsize compiled: \timestamp (local time) --- \textit{SVN-\svnId}
	  }
	%}	
}

% the default definitions of \nbnote and \svnversion are empty	
\newcommand{\nbnote}[2]{}
\newcommand{\svnversion}{}

% the actions perfomed while in comment mode
\newcommand{\commentmode}{
		% \usepackage{draftwatermark} \SetWatermarkLightness{0.70} \SetWatermarkText{\Huge DRAFT}   	
		\renewcommand{\nbnote}[2]{%
		  \fbox{\bfseries\sffamily\scriptsize##1}
		  {\sf\small $\blacktriangleright$ \textit##2 $\blacktriangleleft$}
		  %{\sf\small$\blacktriangleright$\textit{##2}$\blacktriangleleft$}%
		}
	 \renewcommand{\svnversion}{\emph{\footnotesize\svnId}}
	 \AddEverypageHook{\paperref{21cm}{-3mm}{\draftheader}}      
	 \AddEverypageHook{\paperref{21cm}{28.5cm}{\draftheader}}      
}


\newcommand\todo[1]{\nbnote{TO DO}{#1}}

\newcommand{\here}{\nbnote{***}{CONTINUE HERE}}
\newcommand\fix[1]{\nbnote{FIX}{#1}}
\newcommand\jl[1]{\nbnote{JL}{#1}}

\newcommand\rw[1]{\color{blue}#1\color{black}}

% dangerous!
% \DeclareOption{comments}{\commentmode}
% \ProcessOptions 

% LISTINGS  %%%%%%%%%%%%%%%%%%%%%%%%%%%%%%%%%%%%%%%%%%%%%%%%%%%%%%%%%%%%%%%%%%%%%%%%%%%%%%%%%%%%%%%

\usepackage[formats]{listings}
\usepackage[T1]{fontenc}


%\usepackage{pifont}
%\usepackage{winfonts}
%\usepackage{windingbats}
%\usepackage{winpifont}
\usepackage{soul}
%\DisableLigatures{encoding = T1, family = courier }
%\usepackage{beramono}
\usepackage[scaled=.84]{DejaVuSansMono}
%\linespread{1.5}  % mathpazo

%\fontfamily{courier} \selectfont 

\newcommand{\code}[1]{\mbox{\texttt{#1}}} % code font
%\let\code\lstinline
\newcommand{\scode}[1]{\texttt{\small{#1}}} % code font
%\newcommand{\code}[1]{%
%  \lstinline[language=JCop,%
%  					 basicstyle={\ttfamily \footnotesize \selectfont},%
%             keywordstyle={\bfseries \color[rgb]{0,0,0}}]!#1!} % code font
\newcommand{\biglistingsfont}{ \ttfamily \normalsize \selectfont }
\newcommand{\listingsfont}{ \ttfamily \normalsize \selectfont }
\newcommand{\tinylistingsfont}{\rmfamily \fontsize{6.5}{6.5} \selectfont}
% \newcommand{\listingsfont}{ \fontfamily{courier} \scalebox{.5}[1] \scriptsize \selectfont }
%\newcommand{\tinylistingsfont}{\fontfamily{courier} \fontsize{6.5}{5}  \selectfont \scalebox{.5}[1]}

%\newcommand{\listingsfont}{\ttfamily \scriptsize}

% COLORS
% eclipse keywords: \color[rgb]{0.49,0.00,0.33}
% eclipse strings:  \color[rgb]{0.16,0.2,1}
% eclipse comments: \color[rgb]{0.247,0.498,0.372},


% -- JCop  ----------------------------------------------------------------------------------------

\newcommand{\layerclass}{\code{jcop.lang.Layer}\xspace}
\newcommand{\with}{\code{with}\xspace}
\newcommand{\without}{\code{without}\xspace}
\newcommand{\withoutall}{\code{withoutall}\xspace}
\newcommand{\layer}{\code{layer}\xspace}
\newcommand{\before}{\code{before}\xspace}
\newcommand{\after}{\code{after}\xspace}
\newcommand{\finalmod}{\code{final}\xspace}
\newcommand{\abstractmod}{\code{abstract}\xspace}
\newcommand{\proceed}{\code{proceed}\xspace} 
\newcommand{\ind}{\hspace*{1.5em}}



\lstdefinelanguage{ContextJS} {
      classoffset=0,
      keywords={cop, proceed, withLayer, withoutLayer, with, break,%
                layerClass, layerObject, createLayer, %
        case, catch, continue, else, extends,%
                false, finally, function, for, if, var, instanceof,%
                new, %
                return, super, this, throw,%
                true, try, undefined, while, },
      morecomment=[l]//,%    
      morecomment=[s]{/*}{*/},%
      morestring=[b]",%
      % basicstyle=\listingsfont,
      % basicstyle= \ttfamily \scriptsize \selectfont
      keywordstyle = \color[rgb]{0.54,0.27,0.00},
      %commentstyle = \color[rgb]{0.0,0.533,0.8},
      stringstyle = \color[rgb]{0.0,0.533,0.8},
      identifierstyle=,
      showstringspaces=false,
      captionpos=b,
      tabsize=2,
      aboveskip=5mm,
      %numbersep=-10pt,      
      extendedchars=true,
      frame=tb,
      framexleftmargin=0pt,
      framexrightmargin=-2pt,
      numbers=left,
      numberstyle=\tiny    
}


% -- None ---

\lstdefinelanguage{todo} {
  basicstyle=\color{MidnightBlue}\scriptsize\ttfamily,
  breaklines=true,
  xleftmargin=0.5cm,
  extendedchars=true,
  frame=none
}

\lstdefinelanguage{none} {
  keywords={},
  alsoletter={+},
  morekeywords={}
  frame=none
}

% -----------
  
  % general config
  \lstset{
    framexrightmargin=0pt,      
    numberstyle=\tiny,
    numbersep=5pt,      
    frame=tb,
    basicstyle=\linespread{1} \listingsfont, % print whole listing small
    keywordstyle = \bfseries \color[rgb]{0.49,0.00,0.33},%\color[rgb]{0.84,0.27,0.40},
    commentstyle = \color[rgb]{0.247,0.498,0.372},
    stringstyle = \color[rgb]{0.16,0.2,1},    
    captionpos=b,
    tabsize=2,
    aboveskip=5mm,
    showlines=true,
    showstringspaces=false,
    escapechar=\§,
    breakautoindent=false    
  }
  
  
  \newcommand{\lstnolinenumbers}{
  	\lstset{
  	  numbers=none,   	  
      framexleftmargin=4pt,
      xleftmargin=4pt,   
    }
  }
  	
  
  \newcommand{\lstlinenumbers}{ 
  	\lstset{
  	  numbers=left,
      framexleftmargin=15pt,
      xleftmargin=15pt,            
    }
}



% Paperref ==========================================================================
% includes a published note at the top of the title page for local copies

\usepackage[absolute]{textpos}
%\setlength{\TPHorizModule}{0.8\textwidth}
\definecolor{shadecolor}{rgb}{0.8,0.8,0.8} 

\newcommand{\paperref}[3]{%  
  \setlength{\TPHorizModule}{#1}
  \textblockorigin{0mm}{#2}
  \begin{textblock}{1}(0,0)    
      \begin{shaded}#3\end{shaded}    
  \end{textblock}
}



% \usepackage{verbatim}% http://ctan.org/pkg/verbatim
% \newenvironment{note}{\endgraf\color{MidnightBlue}\scriptsize\verbatim}{\endverbatim}

\newenvironment{rewrite}{\endgraf\color{MidnightBlue}}{}


%\usepackage{fancyvrb, xstring}
%\DefineVerbatimEnvironment{note}{Verbatim}{
%    formatcom=\color{MidnightBlue}\scriptsize
%}

%\renewcommand{\FancyVerbFormatLine}[1]{
%   {\StrSubstitute[1]{#1}{☐}{-}}
%}

% \usepackage{listings}

\lstnewenvironment{note}%
  {%
    %\minipage{\textwidth} 
    \lstset{
      basicstyle=\color{MidnightBlue}\scriptsize\ttfamily,
      breaklines=true,
      xleftmargin=0.5cm,
      extendedchars=true,
      frame=none
    }
  }
  {
    %\endminipage
  }

\usepackage{grfext}
\PrependGraphicsExtensions*{.pdf}


% \usepackage[toc]{glossaries}
% \makeglossaries

\usepackage[nopostdot,nonumberlist,nomain,acronym,xindy%,toc
]{glossaries} % nomain, if you define glossaries in a file, and you use \include{INP-00-glossary}
\makeglossaries
\usepackage[xindy]{imakeidx}
\makeindex

\renewcommand*{\glsnamefont}[1]{\textbf{#1}}



%%% Local Variables:
%%% mode: latex
%%% TeX-master: "../Felgentreff_2014_Babelsberg"
%%% End:


% \usepackage{todonotes}

\usepackage{makeidx}
\usepackage{bibentry}

%\iftrue
\iffalse % #TODO
\usepackage{eso-pic}
\usepackage{xcolor}
\newsavebox{\draftPageLine}
  % \newsavebox{\draftWatermark}
  \AddToShipoutPicture{%
  %  \AtPageLowerLeft{\usebox{\draftWatermark}}
    \AtPageUpperLeft{%
      \raisebox{-\height}[\height][0pt]{\usebox{\draftPageLine}}}%
    \AtPageLowerLeft{%
      \raisebox{\depth}[\height][0pt]{\usebox{\draftPageLine}}}%
  }

  \def\draftFont{\rmfamily\bfseries}
  \AtBeginDocument{
    \def\draftInfo{%
      {\draftFont
      Draft Draft Draft%
      \hspace*{3cm}\today\hspace*{3cm}%
      Draft Draft Draft%
      }}
      % Draft revision\gitVtags: \gitAbbrevHash{} %
      % (\gitCommitterIsoDate) \gitCommitterName}
    % \sbox{\draftWatermark}{%
    %   \includegraphics[width=\paperwidth]{img/draft}}
    \sbox{\draftPageLine}{%
      \colorbox{black!10}{%
        % enlarge box vertically by 2/3 lines
        \raisebox{0pt}%
        [\dimexpr .33\baselineskip + \height]%
        [\dimexpr .33\baselineskip + \depth]{%
          \makebox[\paperwidth]{\color{black!50}\draftInfo}}}}
  }
\fi

%
%\marginparsepmm
% \newcommand{\note}[1]{}%{\marginpar{\footnotesize{\textit{#1}}}}

\nobibliography*
%\renewcommand{\baselinestretch}{1}
%\parindent 0cm			%Kein Einzug bei neuen Absatz
%\setlength{\parskip}{2ex}
\setcounter{tocdepth}{1}
\setcounter{secnumdepth}{3}

\addtokomafont{disposition}{\rmfamily}
\addtokomafont{paragraph}{\bfseries}
\setkomafont{pagination}{\normalfont\small}

\renewcommand{\chaptermark}[1]{\markboth{#1}{}}
\renewcommand{\sectionmark}[1]{\markright{#1}{}}
\newcommand{\soabooks}{YourBook\xspace}
\newcounter{pbsavefigurecounter}
% 1: pos [tbh]
% 2: width of the left
% 3: width of the right
% 4: code left
% 5: code right
% 6: caption
% 7: label
%
\newcommand{\listingcaption}[2]{
\vspace{-4mm}
\renewcommand{\figurename}{Listing}
\setcounter{pbsavefigurecounter}{\value{figure}}\setcounter{figure}{\value{lstlisting}}
\caption{#1}
\label{#2}
\setcounter{lstlisting}{\value{figure}}
\setcounter{figure}{\value{pbsavefigurecounter}}}


%\usepackage[amsmath,thmmarks,framed]{ntheorem}
%\theoremstyle{plain}
%\theorembodyfont{\small}

%\usepackage{color}
\definecolor{yellow}{rgb}{1,0.88,0}
%\pagecolor{yellow}
%\renewcommand*\FrameCommand{{\color{gray}\vrule width 5pt \hspace{10pt}}}
%\newframedtheorem{defn}{Definition}

%paul
%\linespread{1.8}
\sloppy  % #TODO begin sloppypar begin sloppypar

%\usepackage{showframe}

\usepackage{tgpagella}       % Majuskelziffern
%\usepackage[osf,sc]{mathpazo} % Minuskelziffern
\linespread{1.05}
\let\sffamily\rmfamily

\usepackage{hyperref}
\hypersetup{colorlinks=true, breaklinks=true, citecolor=black, linkcolor=black, menucolor=black, urlcolor=black}
\newcommand{\CtxMS}{CtxMS\xspace}

% %!TEX root = dissertation.tex
% \chapter*{Glossary}

\newglossaryentry{REPL}{name=REPL, 
	description={Read-Eval-Print Loop}}

\newglossaryentry{AOP}{name=AOP, 
	description={Aspect-oriented Programming}}

\newglossaryentry{COP}{name=COP, 
	description={Context-oriented Programming}}

\newglossaryentry{OOP}{name=OOP, 
	description={Object-oriented Programming}}

\newglossaryentry{IDE}{name=IDE, 
	description={Integrated Development Environment}}

\newglossaryentry{SSDE}{name=SSDE, 
	description={Self-supporting Development Environment}}

\newglossaryentry{HTML}{name=HTML, 
	description={Hypertext Markup Language}}

\newglossaryentry{URL}{name=URL, 
	description={Uniform Resource Locator}}

\newglossaryentry{JSON}{name=JSON, 
	description={JavaScript Object Notation}}

\newglossaryentry{TDD}{name=TDD, 
	description={Test-driven Development}}

\newglossaryentry{UUID}{name=UUID, 
	description={Universally Unique Identifier}}

\newglossaryentry{Morph}{name=Morph, 
	description={Graphical object in Morphic framework}}

\newglossaryentry{Submorph}{name=Submorph, 
	description={Child of a morph in a graphical object composition}}

\newglossaryentry{Tool}{name=Tool, 
	description={Interactive object used to modify other objects}}

\newglossaryentry{Prototype}{name=Prototype, 
	description={Object that serves as a template for the creation of other objects}}

\newglossaryentry{Run-time adaptation}{name=Run-time Adaptation, 
	description={Changing of behavior during execution}}

\newglossaryentry{Part}{name=Part, 
	description={Published morphic object structure that can be (re-)used across worlds}}

\newglossaryentry{Parts Bin}{name=Parts Bin, 
	description={Shared repository of parts}}

\newglossaryentry{Layer}{name=Layer, 
	description={Modularity construct that can be dynamically scoped}}


\newglossaryentry{Web-based}{name=Web-based, 
	description={Software that runs on Web-infrastructure}}

\newglossaryentry{World}{name=World, 
	description={Root of all morphs in the display hierarchy of graphical objects}}

\newglossaryentry{Wiki}{name=Wiki, 
	description={Web-based collaborative authoring environment that permits all users to read, edit, and create Web-content}}

%\newglossaryentry{Development Layer}{name=Development Layer, 
%	description={Layer that is used as scaffolding during development}}

%\newglossaryentry{meta-circular}{name=meta-circular, description={circular dependency of an object and a meta-object}}
%\newglossaryentry{meta-object}{name=meta-object, description={object defines or changes another object}}

% \newglossaryentry{co-evolution}{name=co-evolution, description={}}

%\newglossaryentry{Method}{name=Method, 
%	description={Behavior executed by an object}}

%\newglossaryentry{Message dispatch}{name=message dispatch, 
%	description={Process of finding the method to be executed}}

%\newglossaryentry{Dynamic Extent}{name=dynamic extent, 
%	description={Scope that obeys a stack-like discipline during the execution of a code block}}

%\newglossaryentry{object}{name=object, 
%	description={}}

% \loadglsentries{00_glossary}

%% keine Schusterjungen und Witwen

\widowpenalty=10000
\clubpenalty=10000 

\def\sectionautorefname{Section}
\def\subsectionautorefname{Section}
\def\subsubsectionautorefname{Section}
\def\tableautorefname{Table}
\def\figureautorefname{Figure}

\newcommand\slightspacewide{\kern1em}
\newcommand\slightspace{\kern0.5em}
\newcommand\hairspace{\ifmmode\mskip1mu\else\kern0.08em\fi}
\newcommand{\eg}{e.\hairspace{}g.,\xspace}
\newcommand{\Eg}{E.\hairspace{}g.,\xspace}
\newcommand{\Ie}{I.\hairspace{}e.,\xspace}
\newcommand{\ie}{i.\hairspace{}e.,\xspace}
\newcommand{\cf}[1]{cf.\ {#1}\xspace}

\usepackage{bcprules}
\usepackage{stmaryrd}
\usepackage{amssymb,amsmath,amsthm}
\newtheorem{theorem}{Theorem}
\newtheorem{Corollary}{Corollary}
\newtheorem{Lemma}{Lemma}
\newtheorem{Claim}{Claim}
\newtheorem{Proposition}{Proposition}
\newtheorem{Definition}{Definition}

% \newenvironment{proof}{\noindent\textbf{Proof.~~}}
% {\hfill\ensuremath{\square}\\}

\newcommand{\maybe}{\textbf{\footnotesize{?}}\kern-1pt}

\newcommand{\Euid}{\texttt{\upshape E}}
\newcommand{\Huid}{\texttt{\upshape H}}
\newcommand{\Cuid}{\texttt{\upshape C}}
\newcommand{\Iuid}{\texttt{\upshape I}}

\newcommand{\eval}[3]{\texttt{#1} $\vdash$ \texttt{#2} $\Downarrow$ \texttt{#3}}
\newcommand{\stay}[2]{stay(\texttt{#1}) $=$ \texttt{#2}}
\newcommand{\istay}[2]{i-stay(\texttt{#1}) $=$ \texttt{#2}}
\newcommand{\toSolver}[2]{\texttt{E} $\vdash$ \texttt{#1} $\leadsto$ $#2$}
\newcommand{\fromSolver}[2]{$#1$ $\leadsto$ \texttt{#2}}
\newcommand{\csolve}[2]{$#1$ $\models$ $#2$}
\newcommand{\solve}[2]{\texttt{#1} $\models$ \texttt{#2}}
\newcommand{\cfail}[1]{$\not\models$ $#1$}
\newcommand{\fail}[1]{$\not\models$ \texttt{#1}}
\newcommand{\step}[5]{\texttt{<#1|#2|#3>} $\longrightarrow$
  \texttt{<#4|#5>}}
\newcommand{\tc}[3]{\texttt{#1} $\vdash$ \texttt{#2} : \texttt{#3}}
\newcommand{\tcC}[2]{\texttt{#1} $\vdash$ \texttt{#2}}

\newcommand{\evalH}[4]{\texttt{#1;#2} $\vdash$ \texttt{#3} $\Downarrow$ \texttt{#4}}
\newcommand{\solveH}[3]{\texttt{#1;#2} $\models$ \texttt{#3}}
% \newcommand{\solveH}[3]{\texttt{#1;#2} $\models_{E,H}$ \texttt{#3}}
% \newcommand{\solveH}[3]{\texttt{#1} $\models_{E,H}$ \texttt{#3}}
\newcommand{\failH}[1]{$\not\models_{E,H}$ \texttt{#1}}
\newcommand{\stepH}[9]{\texttt{<#1|#2|#3|#4|#5>} $\longrightarrow$
  \texttt{<#6|#7|#8|#9>}}
\newcommand{\tcH}[4]{\texttt{#1;#2} $\vdash$ \texttt{#3} {\upshape :} \texttt{#4}}
\newcommand{\tcHC}[3]{\texttt{#1;#2} $\vdash$ \texttt{#3}}
\newcommand{\toSolverH}[2]{\texttt{E,H} $\vdash$ \texttt{#1} $\leadsto$ $#2$}
\newcommand{\fromSolverH}[3]{$#1$ $\leadsto$ \texttt{#2,#3}}

\newcommand{\toSolverOO}[6]{\texttt{#1,#2} $\vdash$ \texttt{#3} $\leadsto$ \texttt{#4,#5} $\mid$ $#6$}

\newcommand{\inferrule}[3]{\infrule[#1]{\mbox{#2}}{\mbox{#3}}}
\newcommand{\inferax}[2]{\infax[#1]{\mbox{#2}}}

\newcommand{\judgment}[1]{\begin{flushleft} \fbox{#1}\end{flushleft}}

\newcommand{\receq}{\stackrel{\texttt{r}}{\texttt{=}}}

\newcommand{\lookup}[3]{\textit{lookup}(\texttt{#1},\texttt{#2}) = (\texttt{x$_1$} $\cdots$ \texttt{x$_n$},\texttt{#3})}


\newcommand{\combine}[4]{\texttt{#1;#2} $\vdash$ \texttt{#3} $\leadsto$ \texttt{#4}}
\newcommand{\inline}[9]{\texttt{<#1,#2,#3,#4,#5,#6>} $\leadsto$ \texttt{<#7,#8,#9>}}


\newcommand\enter[9]{%
\def\argA{#1}\def\argB{#2}\def\argC{#3}\def\argD{#4}\def\argE{#5}%
\def\argF{#6}\def\argG{#7}\def\argH{#8}\def\argI{#9}%
\entercontinued
}
\newcommand\entercontinued[2]{%
  \textit{enter}(\texttt{\argA},\texttt{\argB},\texttt{\argC},\texttt{\argD},\texttt{
    \argE},\texttt{\argF},\texttt{x$_1$} $\cdots$ \texttt{x$_n$},\texttt{e$_1$} $\cdots$
  \texttt{e$_n$}) = (\texttt{\argG},\texttt{\argH},\texttt{\argI},\texttt{#1},\texttt{#2})
}

\newcommand\stepHO[9]{%
\def\argA{#1}\def\argB{#2}\def\argC{#3}\def\argD{#4}\def\argE{#5}%
\def\argF{#6}\def\argG{#7}\def\argH{#8}\def\argI{#9}%
\stepHOcontinued
}
\newcommand{\stepHOcontinued}[2]{%
  \texttt{<\argA|\argB|\argC|\argD|\argE|\argF>} $\longrightarrow$ \texttt{<\argG|\argH|\argI|#1|#2>}}


\newcommand\evalBig[9]{%
\def\argA{#1}\def\argB{#2}\def\argC{#3}\def\argD{#4}\def\argE{#5}%
\def\argF{#6}\def\argG{#7}\def\argH{#8}\def\argI{#9}%
\evalBigcontinued
}
\newcommand{\evalBigcontinued}[2]{%
  \texttt{<\argA|\argB|\argC|\argD|\argE|\argF>} $\Downarrow$ \texttt{<\argG|\argH|\argI|#1|#2>}}

\newcommand{\resolve}[4]{\texttt{#3;#4} $\models_\texttt{#1}$ \texttt{#2}}

\newcommand{\Scope}{\texttt{\upshape S}\xspace}
\newcommand{\E}{$\mathbb{E}$}
\newcommand{\Ccal}{$\mathbb{C}$}
\newcommand{\Ical}{$\mathbb{I}$}
\renewcommand{\H}{$\mathbb{H}$}
\newcommand{\reinline}[6]{\texttt{<#1,#2,#3,#4>} $\leadsto$ \texttt{<#5,#6>}}

\newcounter{examples}
\lstnewenvironment{example}[1][]%
{\lstset{%
    basicstyle=\ttfamily,%
    frame=l,%
    #1}}{}

\ifx\C\undefined
\newcommand{\C}{\Ccal}
\else
\renewcommand{\C}{\Ccal}
\fi
\newcommand{\I}{\Ical}

\title{Babelsberg}
\author{Tim Felgentreff}

\begin{document}
% 
% Automatically included file.
% 
% Should contain: 
% * hyphenation; 
% * macros for, eg, product names;
% * acronyms
% 
\hyphenation{Small-talk}
\hyphenation{ContextJS}
\hyphenation{PartsBin}
\hyphenation{using}
\hyphenation{Webwerkstatt}

\newacronym{api}{API}{application programming interface}
\newacronym{repl}{REPL}{Read-Eval-Print Loop}
\newacronym{aop}{AOP}{Aspect-oriented Programming}
\newacronym{cop}{COP}{Context-oriented Programming}
\newacronym{oop}{OOP}{Object-oriented Programming}
\newacronym{ide}{IDE}{Integrated Development Environment}
\newacronym{ssde}{SSDE}{Self-supporting Development Environment}
\newacronym{html}{HTML}{Hypertext Markup Language}
\newacronym{url}{URL}{Uniform Resource Locator}
\newacronym{json}{JSON}{JavaScript Object Notation}
\newacronym{tdd}{TDD}{Test-driven Development}
\newacronym{uuid}{UUID}{Universally Unique Identifier}
\newacronym{ocp}{OCP}{Object-Constraint Programming}
\newacronym{sat}{SAT}{Boolean Satisfiability Problem}
\newacronym{dsl}{DSL}{Domain-specific Language}
\newacronym{jit}{JIT}{Just-in-Time}

\newglossaryentry{Babelsberg}{name=Babelsberg,
  description={Our family of Object-Constraint Programming languages that integrate with existing object-oriented languages}}

\newglossaryentry{ccm}{name={constraint construction mode},
  description={The execution mode that is used to interpret constraint expressions and translate them to the solver}}

\newglossaryentry{z3}{name=Z3,
  description={An SMT solver including theories for reals, integers, booleans, bit-vectors, and records}}

\newglossaryentry{Morph}{name=Morph, 
  description={Graphical object in Morphic framework}}

\newglossaryentry{Submorph}{name=Submorph, 
  description={Child of a morph in a graphical object composition}}

\newglossaryentry{Tool}{name=Tool, 
  description={Interactive object used to modify other objects}}

\newglossaryentry{Prototype}{name=Prototype, 
  description={Object that serves as a template for the creation of other objects}}

\newglossaryentry{Run-time adaptation}{name=Run-time Adaptation, 
  description={Changing of behavior during execution}}

\newglossaryentry{Part}{name=Part, 
  description={Published morphic object structure that can be (re-)used across worlds}}

\newglossaryentry{Parts Bin}{name=Parts Bin, 
  description={Shared repository of parts}}

\newglossaryentry{Layer}{name=Layer, 
  description={Modularity construct that can be dynamically scoped}}

\newglossaryentry{Web-based}{name=Web-based, 
  description={Software that runs on Web-infrastructure}}

\newglossaryentry{World}{name=World, 
  description={Root of all morphs in the display hierarchy of graphical objects}}

\newglossaryentry{Wiki}{name=Wiki, 
  description={Web-based collaborative authoring environment that permits all users to read, edit, and create Web-content}}

%%% Local Variables: 
%%% mode: latex
%%% TeX-master: "Felgentreff_2014_Babelsberg"
%%% End: 

%!TEX root = ../dissertation.tex
\begin{titlepage}
  \enlargethispage{2\baselineskip}
  \usekomafont{title}
  \vspace*{-3.5\baselineskip}
  \vtop{
    \includegraphics[height=2.4cm]{front+back/img/logo}
    \hfill
    \includegraphics[height=2.0cm]{front+back/img/hpi_logo_gr}
  }
  \begin{center}
    \vspace{2\baselineskip}
    % \huge{Lively Webwerkstatt - A Self-sustaining Web-based Authoring Environment}\\
    %% \huge{Lively Webwerkstatt - Evolving A Self-sustaining Web-based Development Environment}\\
    %% Evolving Tools in a Self-supporting Web-based (Collaborative) Development Environment \\
    %% 
    %% 
    %% Webwerkstatt -- Collaboratively Evolving a Self-supporting Prototyping Environment
    {\huge %
    Evolving Tools in a Collaborative Self-supporting Development Environment \\
    %% \huge{Evolving Tools in a Web-based Self-sustaining Development Environment} \\
    %(DRAFT)%
    \par
    }
    \vspace{2\baselineskip}

    {\large
      Dissertation\\ 
      zur Erlangung des akademischen Grades eines \\
      Doktors der Ingenieurwissenschaften \\
      (Doktor-Ingenieur)
      % Doktors der Naturwissenschaften \\
      % (Doctor rerum naturalium)\\ 
      \par


      vorgelegt\\
      der Mathematisch-Naturwissenschaftlichen Fakultät\\
      der Universität Potsdam \\
      %%
      \par

      \vspace{\baselineskip}

      von \\ 
      Dipl.-Ing.\,Jens Lincke\\
      %% 
      \par

      \vspace{\baselineskip}

      betreut von\\
      Prof.\,Dr.\,Robert Hirschfeld\\
      Fachgebiet Software-Architekturen\\
      Hasso-Plattner-Institut für Softwaresystemtechnik\\
      Universität Potsdam\\
      \par

      \vfill
      September 2014
    }
    \end{center}
\end{titlepage}
\cleardoublepage
\pagenumbering{roman}	
\cleardoublepage
\vfill {\centering \textbf{Erklärung}} 

\vspace{3ex}

{\centering Hiermit erkläre ich an Eides statt, dass ich die vorliegende Dissertation selbst angefertigt und nur die im Literaturverzeichnis aufgeführten Quellen und Hilfsmittel verwendet habe.}

\vspace{2ex}

{\centering Diese Dissertation oder Teile davon wurden nicht als Prüfungsarbeit für eine staatliche oder andere wissenschaftliche Prüfung eingereicht. }

\vspace{2ex}

{\centering Ich versichere weiterhin, dass ich diese Arbeit oder eine andere Abhandlung nicht bei einer anderen Fakultät oder einer anderen Universität eingereicht habe}

\vspace{9ex}

{\centering Potsdam, den 29. September 2014}

\vspace{9ex}

{\centering Jens Lincke}




{
\pagestyle{plain}
\markboth{}{}
\chapter*{Abstract}
%!TEX root = abstract_scientific.tex

%%% #Domain
Adapting and developing software tools, to fit them to the task at hand while being used, is the domain of end-user development (EUD). Unlike EUD, the programming tools of a general-purpose development environment cannot usually be adapted in a live and interactive way.  
However, there are exceptions. The combined development and run-time environments of languages, such as Lisp and Smalltalk, allow modifying applications, but also adapting the development environments themselves at run-time. This \emph{self-supporting development} approach shortens feedback loops and makes the development style interactive and explorative.
\emph{Lively Kernel} is a Smalltalk-like self-supporting development environment in the Web-browser, which allows Web-applications to be developed in such a lively way. 
Since the environment is shared, the feedback cycles among collaborating users are also shortened. 
Therefore, a Lively Kernel-based wiki allows the development, sharing, and reuse of each other's creations in an environment which evolves while being used.

%%% #Problem --- #MetaCircularDependencies  
Because of computational reflection and meta-circularity, developing in a self-supporting environment inherently has the danger of breaking one's own tools, which is more severe in a shared environment since changes can and often do affect more than a single user.
Therefore, developers should have the means to change their system at run-time, to try out changes, and share new features or tool adaptations with others in a controlled way. 

%% #Approach --- #LivelyPart and #DevelopmentLayer 
To address this,  we developed two approaches: first, we propose to develop tools as run-time-modifiable \emph{parts} that can be cloned to safely change them. Adapted tools can then be shared via a \emph{parts bin}, so that they can be collaboratively developed. Second, we propose to modularize changes to the base system via \emph{layers} that can be scoped, depending on the execution context, to make the development of the base system safer at run-time by reducing the risk of tools breaking themselves. Furthermore, layers can be used to share such changes in a wiki-like collaboration setting.

%% #Implementation ---  #DeepCloning and #ContextJS
We implemented our approach in \emph{Lively Webwerkstatt}, a collaborative, self-supporting development environment based on Lively Kernel. 
Its programming tools are developed as malleable \emph{parts} that are shared via a \emph{parts bin}. 
\emph{Deep cloning} makes evolving tools safe and also allows alternative ideas to be tried out side by side in the running system.
\emph{Development layers} are implemented using \emph{ContextJS}, our context-oriented JavaScript language extension, which supports domain-specific scoping strategies, that allow restraining behavioral adaptations not only to the dynamic extent of an execution, but also structurally to object composition hierarchies. 

%% Evaluation
We present example artifacts and analyze the repository data of Webwerkstatt, in order to discuss and evaluate our approach. Webwerkstatt has been actively used for more than three years now---not only by a small group of core developers, but also by external users, including our students, for whom it has served as a  shared development environment. During this time, its users successfully worked on their projects, adapted tools, and also helped to evolve Lively Webwerkstatt itself. 











}

\chapter*{Acknowledgements}


{\centering First and foremost, I want to thank Robert Hirschfeld for convincing me to pursue a PhD! He continuously provided the right balance between pushing, supporting, and giving me the freedom to explore, so that I had the chance to learn so many things on so many levels. }

%% Alphabetical ?!? 

\vspace{1ex}

{\centering I also want to thank my colleagues and friends  Malte Appeltauer, Robert Krahn,  Tim Felgentreff, Marcel Taeumel,  Bert Freudenberg,  Bastian Steinert, Carl Friedrich Bolz, Michael Perscheid, and Tobias Pape. Working with them was always inspiring, fruitful, and fun and made the last years here at the Hasso-Plattner-Institute fly bye too fast. }

\vspace{1ex}

{\centering I am thankful to Dan Ingalls for introducing Robert Krahn and me to Lively Kernel and providing us with a seemingly never-ending source of inspiration that lit our fires.  }

\vspace{1ex}

{\centering I also want to thank our graduate students, who helped evolve Lively Webwerkstatt: Julius Dannert, Marko Roeder, Lauritz Thamsen, Fabio Bornhofen, Christopher Schuster, and Astrid Thomschke. }
% I also want to thank the many students of the software design and end-user programming seminars.

\vspace{1ex}

{\centering I want to thank my friends Philipp Engelhard, Stephanie Platz, Anna Filippova, Jenny Rackwitz, Patricia Gerasch, and Philipp Tessenow  for tolerating my continuous explanations of various aspects of my thesis while still giving me valuable feedback from so many perspectives. }

\vspace{1ex}

{\centering I thank my family for their continuous support and encouragement even though there where other things that occupied our minds over the last year. }

\vspace{1ex}

{\centering I would also like to acknowledge the financial support of the Hasso Plattner Design Thinking Research Program for our project.}


{
\tableofcontents
\enlargethispage{\baselineskip}
}
\cleardoublepage
\pagenumbering{arabic}
\cleardoublepage


{
\listoffigures
% \listoftables
% \lstlistoflistings
\glsaddall
%\printglossary[title=List of Abbreviations,toctitle=Abbreviations]
\printglossary[title=List of Technical Terms,toctitle=Terms and Abbreviations]

}

\part{Solving Constraints on Object Behavior}

\chapter{Introduction}
\label{cha:introduction}

\chapter{Problem Statement}
\label{cha:problem-statement}

\section{State of the Art}
\label{sec:state-art}

\section{Motivation and Goals}
\label{sec:design-goals}

\chapter{Relevant Constraint Theory}
\label{sec:relev-constr-theory}

\section{Constraint Solving Strategies}
\label{sec:constr-solv-strat}

\section{Constraint Hierarchies}
\label{sec:constr-hier}

\section{Incremental Re-Satisfaction of Constraints}
\label{sec:incr-re-satisf}

\part{Object-Constraint Programming}
\label{part:object-constr-progr}

\chapter{Design}
\label{sec:design}

\section{Executable Specifications}
\label{sec:exec-spec}

\chapter{Semantics}
\label{sec:semantics}

\section{Constraints over Primitive Types}
\label{sec:primitive-types}

\section{Constraints over Records with Identity}
\label{sec:records-with-ident}

\section{Constraints over Objects and Messages}
\label{sec:constr-with-mess}

\part{Implementations of Object-Constraint Programming}
\label{part:impl-object-constr}

\chapter{Babelsberg/R -- Object-Constraint Programming with VM Support}
\label{part:babelsb-object-const}

\section{Goals for Babelsberg/R}
\label{sec:goals-babelsbergr}

Performance, Full Integration

\section{Implementation}
\label{sec:implementation}

\chapter{Object-Constraint Programming as a Library}
\label{part:object-constr-progr-1}

\section{Goals}
\label{sec:goals}
- Avoid limitations of in-VM approach

- Allow richer tools

\section{Implementation}
\label{sec:implementation-1}

\subsection{Implementation Design}
\label{sec:impl-design}
Host Language requirements, intercession, in-language interpreter...

\subsection{Babelsberg/JS}
\label{sec:babelsbergjs}

\subsection{Babelsberg/S}
\label{sec:babelsbergs}

\chapter{Comparison of Implementation Approaches}
\label{cha:comp-impl-appr}

\section{Conformance to the Semantics}
\label{sec:conf-semant}

\section{Limitations}
\label{sec:lessons-learned}

\subsection{Performance}
\label{sec:perf-eval}

\part{Applications with Object-Constraint Programming}
\label{part:pract-object-constr}

\chapter{Patterns}
\label{sec:patt-object-constr}

\chapter{Control Structures}
\label{sec:contr-struct-object}

\chapter{Tools}
\label{sec:tools}

Debugger, Inspector

\part{Discussion}
\label{part:disc-contr}

\chapter{Related Work}
\label{sec:hist-constr-progr}

\section{Constraint Programming Systems}
\label{sec:constr-progr-syst}

\section{Constraint Logic Programming}
\label{sec:logic-constr-progr}

\section{Constraint Imperative Programming}
\label{sec:constr-imper-progr}

\section{Constraint Programming Libraries and DSLs}
\label{sec:constr-progr-libr}

\section{Current Approaches for Constraints in General Purpose Programming}
\label{sec:curr-appr-constr}

\chapter{Discussion}
\label{cha:discussion}

\chapter{Future Work}
\label{cha:future-work}

\chapter{Conclusions}
\label{cha:conclusions}



%%% Local Variables:
%%% mode: latex
%%% TeX-master: "Felgentreff_2014_Babelsberg"
%%% End:

\chapter*{Publications}

\section*{Journals}

\begin{itemize}
\item \fullcite{FelgentreffEST15}
\item \fullcite{HirschfeldJSSST15}
\item \fullcite{FelgentreffJOT14}
\end{itemize}

\section*{Conferences}

\begin{itemize}
\item \fullcite{FelgentreffOOPSLA2015}
\item \fullcite{FelgentreffDLS2015}
\item \fullcite{HirschfeldJSSST14}
\item \fullcite{SteinertDLS2014}
\item \fullcite{FreudenbergDLS2014}
\item \fullcite{FelgentreffECOOP14}
\item \fullcite{PerscheidWCRE2014}
\end{itemize}

\section*{Workshops}

\begin{itemize}
\item \fullcite{FelgentreffFPW2015}
\item \fullcite{LehmannCOP2015}
\item \fullcite{NiephausICOOOLPS2015}
\item \fullcite{FelgentreffICOOOLPS2015}
\item \fullcite{GraberREBLS2014}
\item \fullcite{TaeumelCOP2014}
\item \fullcite{FelgentreffWASDETT2013}
\end{itemize}

% \vbox{
% \section*{Book Chapters}

% \begin{itemize}
% 	\item \fullcite{BookEntry}
% \end{itemize}
% }

\section*{Technical Reports}
\begin{itemize}
\item \fullcite{felgentreff:2014:semantics}
\item \fullcite{felgentreff+:2014:babelsberg}
\end{itemize}

%%% Local Variables:
%%% mode: latex
%%% TeX-master: "Felgentreff_2015_Babelsberg"
%%% End:

\bibliographystyle{plain} % alpha
\bibliography{thesis}
\end{document}

%%% Local Variables:
%%% mode: latex
%%% TeX-master: t
%%% End:
